% SHARD-ID: RC-04V-DENINGER-CUSP
% SHARD-TITLE: The Phantasm XII: the Hecke Correspondence with a Cusp. The S-Arithmetic Quotient as a Graph of Modular Surfaces, the Eisenstein Classes Are Neither the Cup Radical Nor the Poles Nor the Exit, the Prime-Level Scattering Matrix of Gamma_1(q) with All Constants, the Twisted Theta Edges, the Pure Blaschke Character Symbols, and the Level Model with Its Local Factors
% SHARD-SUMMARY: Gamma_1(N)[1/p] \ (H x T_{p+1}) is a graph of surfaces with one edge (two vertex slices Y, one edge slice Y_Delta, degeneracy maps of degree p + 1), not a two-dimensional square complex; the vertex slice's H^1(Y_1(N)) splits Hecke-equivariantly (Eichler--Shimura, Manin--Drinfeld) as parabolic plus Eisenstein, with Eisenstein eigenvalues chi_1(p) + p chi_2(p) (p + 1 only for the trivial packet); the ordinary cup product on H^1(Y) is zero, the compact-support pairing is perfect, and the four-fold identity Eisenstein = radical = even poles = exit is refuted in dimension and parity, while the Eisenstein root moduli 1, p exclude any positive p-similitude metric there.
% SHARD-SUMMARY: For Gamma_1(q), q prime, the scattering matrix in the diamond-Fourier basis has, per even nebentypus psi, the block [[0, A_psi],[A_{bar psi}, 0]] with A_psi(s) = q^{1/2-s} Lambda_psi(2s-1)/Lambda_psi(2s), and the trivial block phi(s)/(q^{2s}-1) [[q-1, q^s - q^{1-s}],[q^s - q^{1-s}, q-1]] with rank-one residue 3/(pi(q+1)); squarefree trivial level is a tensor product of such blocks; the determinant is explicit. The level-scaled twisted Siegel theta has its two edges at the two cusps (theta_chi(t/y) and sqrt(y/t) theta_chi(fty)), giving A_chi as a quotient of two GL_1 theta channels; odd characters vanish in weight zero; the conjugation-correct symbol Lambda_chi(1+2i tau)/Lambda_{bar chi}(1-2i tau) is pure Blaschke unconditionally.
% SHARD-SUMMARY: The prime-level inverse symbol factors as arithmetic (Theta_chi) times local inner factors u b_{-+a}(u) (u = q^{i tau}, a = q^{-1/2}) and delays; K_AB = K_A + A K_B; the arithmetic reduction is two copies of K_{Theta_chi} per even character; the exit rank is exactly h = q - 1 (one exit per scalar coordinate; finite-time defects have infinite rank); the spectral positivity criterion is GRH for the represented characters plus no Jordan chains; no boundedly equivalent positive metric exists. asm:h-arith-gamma1 is refuted as a Q statement.
% SHARD-KEYWORDS: Ihara lattice, graph of spaces, Hecke correspondence, Eichler--Shimura, Manin--Drinfeld, mixed Hodge structure, Eisenstein cohomology, cup product, compact support, scattering matrix, Gamma_1(N), cusps, nebentypus, Dirichlet L-function, Gauss sum, twisted theta, Blaschke product, delay, exit rank, GRH
% SHARD-NOTES: none
% SHARD-SCRIPTS: scripts/deninger_cusp.py

\section{The Phantasm XII: the Hecke Correspondence with a Cusp}
\label{sec:deninger-cusp}

Round of 2026-09-22 (brief \texttt{notes/deninger-cusp/astra-brief.md}, prover \texttt{notes/deninger-cusp/astra-proofs.md},
numerics \texttt{scripts/deninger\_cusp.py}, review \texttt{notes/reviews/deninger-cusp-2026-09-22.md}). The brief proposed
the non-cocompact twin of \cref{sec:deninger-lps}: Ihara's quotient $\Gamma_1(N)[1/p]\backslash(\mathbb H\times T_{p+1})$ with
horizontal cohomology $H^1(Y_1(N))$, transverse transport $T_p$, Deninger's package on the cuspidal part and the
Eisenstein classes as radical, poles and exit at once, and the level-$N$ scattering matrix as a diagonal matrix of GL$_1$
theta-channel edge quotients (H-ARITH over $\mathbb Q$, \cref{asm:h-arith-gamma1}). Most of that was refuted and replaced
by explicit statements. In this shard $q$ is the level prime, $p$ the transverse prime, $u = e^{i\tauM\log q}$.

\subsection{The arithmetic correspondence and what the Eisenstein classes are not}

\begin{proposition}[The $S$-arithmetic quotient is a graph of surfaces; the vertex cohomology and its Eisenstein part]
\label{prop:ihara-graph-of-surfaces}
\claimstatus{prop:ihara-graph-of-surfaces}{sketched}
For $\Gamma_S = \{\gamma\in\mathrm{SL}_2(\mathbb Z[1/p]):\gamma\bmod N\text{ upper unipotent}\}$ (a finite-index subgroup of
Ihara's lattice, not byte-cited), the action on $T_{p+1}$ has quotient one edge with two vertex types, stabilisers
$\Gamma = \Gamma_1(N)$, $g\Gamma g^{-1}$ ($g = \mathrm{diag}(1,p)$) and $\Delta = \Gamma_1(N)\cap\Gamma_0(p)$ (reduction modulo $N$ is
onto by the Chinese remainder theorem, so the level creates no extra orbits); the quotient of $\mathbb H\times T$ is the graph
of surfaces $Y\xleftarrow{\alpha}Y_\Delta\xrightarrow{\beta}Y$, $\alpha(z) = z$, $\beta(z) = pz$, both of degree $p+1$. It is
three-dimensional with two vertex slices, not a square complex; the $p+1$ ``branches'' are the sheets of the degeneracy maps.
A vertex slice has $H^1(Y;\mathbb R)$ of dimension $2g+h-1$; horizontal cohomology of the whole quotient is a coefficient
system over the edge, not this single space. Eichler--Shimura and Manin--Drinfeld (not byte-cited) give the
Hecke-equivariant mixed-Hodge splitting $0\to H^1(X)\to H^1(Y)\to\tilde H^0(C)(-1)\to0$, unique over $\mathbb Q$ (no morphism
from weight two to weight one), $\mathrm{Eis}$ of dimension $h-1$ and type $(1,1)$; the transport
$\mathcal T_p = \alpha_{*}\beta^{*}$ on forms (sum over local sheets) commutes with $d$ and preserves compact supports, and its
cohomology action is the geometric $T_p$ (definition), identified with the holomorphic Hecke operator by
Eichler--Shimura (theorem). Eisenstein eigenpackets with primitive inducing characters $(\chi_1,\chi_2)$ have
$a_p = \chi_1(p)+p\chi_2(p)$ ($p\nmid N$); $p+1$ holds for the trivial pair only. At prime level $\ell\ge5$ there is one trivial
packet and two packets $(1,\chi),(\chi,1)$ per nontrivial even $\chi$, total $\ell-2 = h-1$.
\end{proposition}

\begin{theorem}[Frobenius, the three pairings, and the refuted four-fold identity]
\label{thm:eisenstein-not-radical-not-exit}
\claimstatus{thm:eisenstein-not-radical-not-exit}{sketched}
(a) For a weight-two newform of nebentypus $\psi$ the Frobenius polynomial on its two-dimensional constituent is
$X^2-a_pX+p\psi(p)$ (Eichler--Shimura congruence relation, Igusa; not byte-cited), and on $H^1_{\mathrm{et}}(X)$ Frobenius is
semisimple with all roots of modulus $\sqrt p$ (Weil, Tate) and satisfies $F_p^{t}\Omega_XF_p = p\Omega_X$ for the cup form. % bare-ok
The brief's Bass companion on $H^1(X)\oplus H^1(X)$ has dimension $4g$, not $2g$, omits $\psi(p)$, and is not Frobenius; a
companion can represent one cyclic two-plane. Poincar\'e duality gives the \emph{adjoint} identity
$\Omega_X(T_px,y) = \Omega_X(x,\langle p\rangle^{-1}T_py)$, not a similitude. (b) Bass positivity for a real symmetric $T$ on % bare-ok
$(K,M)$: $G_B = \begin{psmallmatrix}M&-MT/2\\-MT/2&pM\end{psmallmatrix}$ is positive iff $|\lambda|<2\sqrt p$ strictly (equality is
a Jordan block); the doubled certificate, not the undoubled Petersson metric, is the invariant metric. (c) For a
punctured curve $H^2(Y) = 0$, so the ordinary cup product on $H^1(Y)$ is identically zero (radical everything), while
$H^1_c(Y)\times H^1(Y)\to H^2_c(Y)$ is perfect (no radical); the skew form on $H^1_c$ has radical $\ker j = H^0(C)/H^0(X)$ of
dimension $h-1$ and weight zero; the split form $\Omega_{\mathrm{split}}(x,y) = \Omega_X(rx,ry)$ through the Manin--Drinfeld % bare-ok
retraction has radical exactly $\mathrm{Eis}$, but is not the cup product and not a face pairing. (d) Eisenstein packets
have Frobenius roots $\chi_1(p),p\chi_2(p)$ of moduli $1,p$: no positive $p$-similitude metric exists on them (Hermitian
version for complex packets). The four objects have different dimensions and parities: boundary cohomology $h-1$, the
continuous spectrum multiplicity $h$, the residual constant one, $H^0(X)$ and $H^2(X)$ two even lines; $Y(1)$ has $h = 1$,
$H^1 = 0$ and a nonconstant scattering coefficient with a pole at $s = 1$. Hence neither the scattering exit nor the pole
pair is the degree-one Eisenstein radical. The gamma factor and the $T\log T$ zero count live in the completed
$L$-ratios of the continuous families, not in any finite-dimensional cohomology; the LPS model already has even poles
$1,17$ without a cusp and an odd-dimensional horizontal precursor (\cref{sec:deninger-lps}).
\end{theorem}

\subsection{The prime-level scattering matrix with all constants}

\begin{theorem}[Cusps of $\Gamma_1(N)$; the $\Gamma_1(q)$ scattering matrix]
\label{thm:prime-level-scattering}
\claimstatus{thm:prime-level-scattering}{sketched}
(a) The cusps of $\Gamma_1(N)$ ($N\ge5$) are $\{(c\bmod N,a\bmod d): d = (c,N), a\in(\mathbb Z/d)^\times\}/\pm$, a parametrised set
with $h_1(N) = \tfrac12\sum_{d|N}\varphi(d)\varphi(N/d)$ ($4,6,10,12$ at $N = 5,7,11,13$); the brief's $\gcd(d,N/d)$ count is
$\Gamma_0(N)$'s. At prime level $q$ the diamond group $G = (\mathbb Z/q)^\times/\pm$ acts freely and transitively on each of two cusp
orbits, so Fourier transform on $G$ gives a two-dimensional multiplicity space per even nebentypus $\psi$, including real
$\psi$: not the inverse-pair blocks of \cref{prop:cusps-are-class-group-bond}, which concerns $\mathrm{GL}_2(O(C\setminus P_0))$.
(b) With width-one scaling matrices, $E^\psi_\infty(z,s) = \tfrac12\sum_{(c,d) = 1,\ q|c}\bar\psi(d)y^s|cz+d|^{-2s}$,
$E^\psi_0 = E^{\bar\psi}_\infty|W_q$, and $\Phi_\psi(s)_{ab}$ defined by $(E^\psi_b|\sigma_a)_0 = \delta_{ab}y^s+\Phi_\psi(s)_{ab}y^{1-s}$
(rows observing, columns source), for nontrivial even $\psi$ (all primitive):
\[
\Phi_\psi(s) = \begin{pmatrix}0&A_\psi(s)\\A_{\bar\psi}(s)&0\end{pmatrix},\qquad
A_\psi(s) = q^{1/2-s}\frac{\Lambda_\psi(2s-1)}{\Lambda_\psi(2s)} = q^{1/2-s}\frac{\tau(\psi)}{\sqrt q}\frac{\Lambda_{\bar\psi}(2-2s)}{\Lambda_\psi(2s)},
\]
$\det\Phi_\psi = -q^{1-2s}R_\psi R_{\bar\psi}$, $A_\psi(s)A_{\bar\psi}(1-s) = 1$ (the matrix functional equation), unitary on
$\re s = \tfrac12$, regular at $s = 1$ with $A_\psi(1) = \pi L(1,\psi)/(qL(2,\psi))$ (zero residue); the incoming series has no
$y^{1-s}$ term at its own cusp (a nontrivial character averages to zero). For real $\psi$ the fixed vectors $(1,\pm1)$ give
two scalar channels $\pm A_\psi$; for complex $\psi$ no fixed similarity diagonalises the block for all $s$. (c) Trivial
nebentypus: with $\phi(s) = \Lambda_0(2s-1)/\Lambda_0(2s)$ and the oldform matrix $H(s) = \begin{psmallmatrix}1&q^s\\q^s&1\end{psmallmatrix}$,
\[
\Phi_1(s) = \phi(s)H(1-s)H(s)^{-1} = \frac{\phi(s)}{q^{2s}-1}\begin{pmatrix}q-1&q^s-q^{1-s}\\q^s-q^{1-s}&q-1\end{pmatrix},\quad
\phi_\pm = \phi\frac{1\pm q^{1-s}}{1\pm q^s},\quad\det\Phi_1 = \phi^2\frac{1-q^{2-2s}}{1-q^{2s}},
\]
residue at $s = 1$ equal to $\tfrac{3}{\pi(q+1)}$ times the all-ones matrix ($1/\mathrm{vol}(Y_0(q))$ per entry). (d) Squarefree
trivial level: $\Phi_{0,N} = \phi\bigotimes_{r|N}\tfrac1{r^{2s}-1}\begin{psmallmatrix}r-1&r^s-r^{1-s}\\r^s-r^{1-s}&r-1\end{psmallmatrix}$,
determinant $\phi^h\prod_{r|N}((1-r^{2-2s})/(1-r^{2s}))^{h/2}$, $h = 2^{\omega(N)}$: no Dirichlet ratio, the scope of
\cref{obs:level-tower-correction} with ``squarefree'' retained. (e) $\det\Phi_{\Gamma_1(q)} = \phi^2\tfrac{1-q^{2-2s}}{1-q^{2s}}
(-1)^{m-1}q^{(1-2s)(m-1)}\prod_{\psi\ne1\text{ even}}R_\psi^2$, $m = (q-1)/2$: each $L$-ratio squared; the residue in the width-one
basis is the all-ones matrix times $6/(\pi(q^2-1))$. At general level the Eisenstein space is spanned by primitive inducing
pairs $(\chi_1,\chi_2,b)$ and their level raises (Huxley 1984, Hejhal 1983; explicit change of basis in Young 2017; not
byte-cited); only the spanning structure is used. Verified: \cref{num:deninger-cusp}.
\end{theorem}

\begin{theorem}[Twisted theta edges; the character symbol is pure Blaschke]
\label{thm:twisted-theta-edges}
\claimstatus{thm:twisted-theta-edges}{sketched}
For a primitive even $\chi$ of conductor $f$, $\theta_\chi(v) = \sum_n\chi(n)e^{-\pi n^2v/f}$ and the level-scaled lattice sum
$\mathcal T_\chi(z,t) = \sum_{m,n}\chi(n)e^{-\pi t|fmz+n|^2/(fy)}$ (periodic in $x$): $\tfrac12\int_0^\infty\mathcal T_\chi t^s\,d^\times t
= \Lambda_\chi(2s)E^{\bar\chi}_\infty(z,s)$, and the constant terms are $\int_0^1\mathcal T_\chi(x+iy,t)dx = \theta_\chi(t/y)$ at $\infty$
and $\sqrt{y/t}\,\theta_\chi(fty)$ at $0$: the two edges are at the two cusps (a nontrivial primitive twist has a zero
character average, so the brief's two terms at one cusp are false), with Mellin edges $\Lambda_\chi(2s)y^s$ and
$f^{1/2-s}\Lambda_\chi(2s-1)y^{1-s}$, whose quotient is $A_\chi(s)$ of \cref{thm:prime-level-scattering}: a GL$_1$ theta-edge
derivation of the scattering entry. Poisson with constants: $\theta_\chi(v) = \epsilon_\chi v^{-1/2}\theta_{\bar\chi}(1/v)$,
the weighted involution $\mathcal WF(t) = t^{-1}F(1/t)$ exchanges the symmetrically scaled edges with $\chi\to\bar\chi$,
$\tfrac12\int A_\chi t^s\,d^\times t = \Lambda_\chi(2s)y^s$, $\tfrac12\int B_\chi t^s\,d^\times t = \Lambda_\chi(2s-1)y^{1-s}$ for all
$s$; the functional equation exchanges the character \emph{and} the argument. Odd $\chi$: the unsigned Gaussian sums vanish
identically and there is no weight-zero scattering matrix (in particular none for the order-four character modulo $5$); the
GL$_1$ replacement is $\theta_{\chi,a}(v) = \sum n^a\chi(n)e^{-\pi n^2v/f}$ with $\tfrac12\int\theta_{\chi,a}v^{(u+a)/2}d^\times v =
\Lambda_\chi(u)$. The conjugation-correct symbol $E_\chi(\tauM) = \Lambda_{\bar\chi}(1-2i\tauM)$,
$\Theta_\chi = E_\chi^\#/E_\chi = \Lambda_\chi(1+2i\tauM)/\Lambda_{\bar\chi}(1-2i\tauM) = \epsilon_\chi\Lambda_\chi(1+2i\tauM)/\Lambda_\chi(2i\tauM)$
(the same-character quotient is not unimodular for complex $\chi$), with $A_\chi(\tfrac12+i\tauM) = f^{-i\tauM}\epsilon_\chi\Theta_\chi^{-1}$,
is inner and pure Blaschke unconditionally: zeros $w_\zero = \gamma/2+i(1-\sigs)/2$ in $\mathbb C_+$ with multiplicity, no
singular factor, no delay ($\Theta_\chi(iv)\sim\epsilon_\chi\sqrt{\pi/(fv)}$), by the real entire $\epsilon_\chi^{-1/2}\Lambda_\chi(\tfrac12+iz)$
and its genus-one Hadamard product, as in \cref{thm:symbol-edge-quotient} for $\chi = 1$ (where the pole-removing $r$ of
\cref{prop:siegel-constant-term} appears; Uetake's reduced causal factor is $\cxi(2s)/\cxi(-2s)$,
\texttt{uetake-2007:paper.txt:615-680}).
\end{theorem}

\subsection{The level model: arithmetic modes, local factors, delays, exits}

\begin{theorem}[The prime-level inner model and its decomposition]
\label{thm:level-model-decomposition}
\claimstatus{thm:level-model-decomposition}{sketched}
Let $a = q^{-1/2}$, $b_c(u) = (u-c)/(1-cu)$, $L_+ = ub_{-a}(u)$, $L_- = ub_a(u)$, $r(\tauM) = (\tauM-i/2)/(\tauM+i/2)$. Using the
inverse of the Eisenstein coefficient matrix with its bound-state pole removed (\cref{sec:cusp-graph-scattering},
\cref{sec:h-theta-channels}), $\phi_\pm(\tfrac12+i\tauM)^{-1} = r^{-1}\Theta_1L_\pm$, and the two trivial-block inner entries are
$I_+ = \Theta_1L_+$, $I_- = \Theta_1L_-/r$; each nontrivial even $\psi$ gives $I_{\psi,1} = e^{i\tauM\log q}\Theta_\psi$,
$I_{\psi,2} = e^{i\tauM\log q}\Theta_{\bar\psi}$ up to constant phases and a right permutation. Hence
$K^{\mathrm{scat}}_q\simeq K_{I_+}\oplus K_{I_-}\oplus\bigoplus_{\psi\ne1}(K_{e^{i\tauM\log q}\Theta_\psi}\oplus K_{e^{i\tauM\log q}\Theta_{\bar\psi}})$,
including local factors and delays; since $K_{AB} = K_A\oplus AK_B$, discarding them defines a separate \emph{arithmetic
reduction} $K^{\mathrm{arith}}_q\simeq\bigoplus_{\chi\text{ even}}K_{\Theta_\chi}^{\oplus2}$: two copies per even character
(also for $\chi = 1$ and quadratic $\chi$), not one per character and no odd characters; the delay $e^{i\tauM\log q}$ has model
space $L^2(0,\log q)$ with no zeros, and the local factors have periodic zeros at height $\tfrac12$ (the zero of $L_-$ at
$i/2$ is the cancelled one). At $q = 5$: four cusp coordinates, two zeta channels and two quadratic-character channels.
Exit: for a matrix inner $\Theta$, $I-Z^{*}Z = J^{*}J$ with $J^{*}e = (\Theta(w)-\Theta(0))e/w$, $JJ^{*} = 1-\Theta(0)^{*}\Theta(0)$, so
$\mathrm{rank}\,J = \mathrm{rank}(1-\Theta(0)^{*}\Theta(0))\le h$ with equality iff no constant isometric direction; for the
prime-level symbols every entry is nonconstant inner, so $\mathrm{rank}\,J = q-1 = h$ exactly: one exit per scalar coordinate,
not per character or per zero. In continuous time the generator has one boundary exit per coordinate, but the finite-time
defect $I-C_\tsg^{*}C_\tsg$ has infinite rank (independent restrictions of infinitely many exponentials to $[0,\tsg]$): the
disc shift is a cogenerator convention. Modes: $\mathcal M^{-1}K_{\Theta_\chi} = \overline{\mathrm{span}}\{(\log y)^jy^{-(1-\sigs)/2-i\gamma/2}1_{y>1}\}$
over the zeros of $L(s,\chi)$ with jets, complete by divisibility, damped uniformly by $y^{-1/4}$ iff GRH for $\chi$; a
diamond $\langle d\rangle$ acts on the whole $\psi$-block by $\psi(d)$ while the $L$-labels are $\psi,\bar\psi$ (in general the
diamond label is $\chi_1\chi_2$ and the $L$-character $\chi_1\bar\chi_2$): ``$L$-character $=$ diamond character, one channel
each'' is refuted. Only the principal completed $L$ has poles; the pole-even/zero-odd grading of $L$-functions does not
make the Eisenstein classes even.
\end{theorem}

\begin{proposition}[Positivity on the level model: GRH plus no Jordan chains; no boundedly equivalent metric]
\label{prop:level-positivity-criterion}
\claimstatus{prop:level-positivity-criterion}{sketched-conditional}
On a finite, functional-equation-stable selection of arithmetic zero modes with complete chains and partners, with
$D = A+\tfrac14$ and eigenvalues $d_\zero = (\sigs-\tfrac12)/2-i\gamma/2$, a positive metric making $e^{\tsg D}$ unitary for all real
$\tsg$ exists iff all selected zeros are critical and every Jordan chain has length one (repeated simple zeros in
independent cusp coordinates are allowed; a multiple analytic zero in one scalar factor is not): the infinitesimal form
of the finite polarisation criterion, ``semisimple with critical spectrum'', not ``simple spectrum''. A pairing pairs
$(\chi,\zero)$ with $(\bar\chi,1-\zero)$. On the algebraic direct sum of arithmetic modes a positive diagonal completion
exists iff GRH and simplicity for all \emph{represented} $L$-functions (even characters at prime level; not all
characters modulo $N$); it does not apply to the full model with local modes and delays, and proves neither GRH nor
simplicity. Under GRH and simplicity for one represented character, two normalised kernels at height $\tfrac14$ have
overlap $\tfrac12/\sqrt{(a-a')^2+\tfrac14}$, the $T\log T$ count forces gaps tending to zero, so no bounded, boundedly
invertible positive metric makes the centred evolution unitary (the character extension of
\cref{prop:kernels-not-riesz} and \cref{thm:no-energy-equivalent-metric}); a spectral completion making all modes
orthonormal changes the exit to $(1-e^{-\tsg/2})I$, of infinite rank.
\end{proposition}

\begin{observation}[H-ARITH over $\mathbb Q$: the assumption is refuted as stated; what is open]
\label{obs:h-arith-q-verdict}
\claimstatus{obs:h-arith-q-verdict}{sketched}
The literal $\mathbb Q$-analogue of \cref{asm:h-arith-gamma1} (one diagonal monomial-times-$L$ channel per Dirichlet character,
diamonds acting by the $L$-character) is refuted by \cref{thm:prime-level-scattering,thm:level-model-decomposition}:
two cusp coordinates per even nebentypus, rational local factors already in the trivial block, odd characters absent
in weight zero, and diamond label $\ne$ $L$-label. What survives: per even character the completed $L$-ratio with the
conductor power. Over $\mathbb F_q[T]$ the double-coset, inducing-pair, constant-term and finite-Fourier structure has
analogues, while the Gaussian, the gamma factor, the continuous clock and the $T\log T$ argument do not transfer; a
cubic-level calculation must fix the subgroup (determinant and centre), singular cusps, heights, local bases and the
infinity representation (lane \texttt{notes/h-arith-fqt/}). $\det S = (-1)^hp/\tilde p$ of \cref{thm:multi-exit-scattering} is
rational in the tree variable; the archimedean matrices have completed $L$-functions, infinitely many modes and
delays, so no finite $p$ represents them, and determinant agreement loses directional cancellations. The pole pair of
the completed zeta is not the Eisenstein radical of $H^1$ of this quotient; $\cxi$ is entire.
\end{observation}

\begin{numerical}[Numerics lane, the cusp correspondence and level $N$]
\label{num:deninger-cusp}
\claimstatus{num:deninger-cusp}{numerical}
\texttt{scripts/deninger\_cusp.py} ($163$ assertions, $115$ s, mpmath; \texttt{outputs/deninger\_cusp.txt}; ledger in
\texttt{notes/deninger-cusp/numerics.md}; $163$ pass; written without reading the prover's scripts): the six table entries
of the $\psi$-blocks for $\chi_5,\chi_7,\chi_{13}$ at two points to $12$ decimals, the Gauss-sum form to $10^{-20}$,
$\Phi_\psi(s)\Phi_\psi(1-s) = I$ and unitarity on the critical line to $10^{-25}$, the determinants, the trivial block, its
eigenvalues and residue at $q = 5,7$, the full $4\times4$ $\Gamma_1(5)$ matrix (functional equation, residue every entry $1/(4\pi)$
of rank one), the $\Gamma_1(q)$ determinant, $A_\psi(1) = 0.38293620316254140264$ for $(\cdot/5)$; Poisson for four characters by
direct summation, the $A_\chi,B_\chi,C_\chi$ table at $y = 1.3$, $t = 0.7$, the constant terms at both cusps by two-dimensional
lattice sums (errors $\le4\cdot10^{-16}$), the odd unsigned sum identically zero, the Mellin edges and their quotient $= A_\chi$,
$|\Theta_\chi| = 1$ on $\mathbb R$ and $<1$ on an upper-half-plane grid, the delay identity, the asymptotic ratio $1.0031$ at
$v = 40$, the non-unimodularity ($0.68$) of the same-character quotient for $\chi_7$; the zeros of $L(s,(\cdot/5))$
$\gam{1} = 6.648453344727714716$, $\gam{2} = 9.831444432886669616$ to $21$ digits, $\Theta(w_j) = 0$, the mode values, the kernel
transform and the overlap $0.2997$; the Eisenstein multiset $\{p+1\}\cup\{\chi_1(p)+p\chi_2(p)\}$ for $N = 11$, $p = 2,3$
against the prover's polynomial to $5\cdot10^{-22}$. Findings: only the first form of $A_\psi$ makes the odd diagnostic
meaningful; \texttt{mp.quad} on $[0,\infty)$ needs a subdivided range for the oscillatory mode transform.
\end{numerical}

\begin{observation}[What this changes; next]
\label{obs:deninger-cusp-next}
\claimstatus{obs:deninger-cusp-next}{sketched}
The level-$q$ modular surface's channel is arithmetic modes times local inner factors with one exit per cusp coordinate;
the local factors are the notebook's finite cusp-graph cavity of the tree quotient at $q$ (lane
\texttt{notes/level-local-cavity/}); Deninger's package holds on cusp forms through Eichler--Shimura and fails on the
Eisenstein classes for the reason the finite note gives (moduli $1,p$), which is a statement about weight, not about
exits. H-ARITH over $\mathbb F_q[T]$ (lane \texttt{notes/h-arith-fqt/}) is where the whole structure is finite and exact.
\end{observation}
